\section{Manual de implementación: convertir la receta de entrenamiento en un proceso ejecutable para el equipo}

\subsection{Ejemplo de ritmo de iteración de cuatro semanas}
Para que el método de esta clase realmente se lleve a la práctica, lo más útil es convertir la receta de entrenamiento en un ritmo fijo, en lugar de experimentos improvisados. A continuación se presenta una plantilla de ritmo de cuatro semanas, utilizable en equipos de investigación de tamaño mediano:
\begin{enumerate}
    \item Semana 1: congelar la tarea objetivo y el protocolo de evaluación, y completar la descontaminación de datos y la ejecución de la línea base;
    \item Semana 2: experimentos de mezcla de datos de SFT, para confirmar el intervalo de estabilidad del comportamiento;
    \item Semana 3: barridos a pequeña escala de DPO/PPO, para verificar la ganancia de preferencia y los efectos secundarios;
    \item Semana 4: experimento conjunto de RLVR y test-time scaling, con el que se obtiene la curva de costo de despliegue.
\end{enumerate}

\begin{importantbox}{Hacer que el «ritmo experimental» importe más que la «inspiración»}
La mejora de la capacidad de razonamiento casi nunca es un salto único, sino la acumulación de pequeños cambios continuos. Solo al fijar los objetivos semanales, las métricas de evaluación y las pruebas de regresión se puede evitar que el equipo repita el mismo error una y otra vez.
\end{importantbox}

\subsection{Plantilla de registro de experimentos: campos mínimos rastreables}
\begin{table}[H]
\centering
\renewcommand{\arraystretch}{1.2}
\begin{tabular}{p{3.0cm}p{3.7cm}p{4.1cm}}
\toprule
\textbf{Campo} & \textbf{Contenido que se registra} & \textbf{Consecuencia de omitirlo} \\
\midrule
Versión de datos & Origen de las muestras, reglas de deduplicación, proporción de mezcla, estado de licencia & No se puede determinar si la ganancia proviene del algoritmo o de una deriva de los datos \\
Configuración de entrenamiento & Optimizador, tasa de aprendizaje, lote, número de pasos de entrenamiento, seed & Los resultados no son reproducibles y es difícil hacer una regresión \\
Definición de la recompensa & Reglas de preferencia/versión del verificador/umbral & Es difícil diagnosticar el reward hacking \\
Protocolo de evaluación & Conjunto de referencias, reglas de descontaminación, intervalo de confianza & Las puntuaciones no son comparables y las decisiones se distorsionan \\
Presupuesto de despliegue & Latencia promedio, P95, costo de tokens, tasa de reintento por fallas & No se puede determinar si es viable ponerlo en producción \\
\bottomrule
\end{tabular}
\caption{Campos mínimos de registro para institucionalizar la receta de entrenamiento}
\end{table}

\begin{knowledgebox}{El punto que más se pasa por alto en la ingeniería}
Muchos equipos solo registran el «mejor resultado» y no registran los «experimentos fallidos». Pero en el entrenamiento de reasoning, las trayectorias fallidas suelen revelar mejor que las trayectorias exitosas problemas de contaminación de datos, sesgo de recompensa o de la estrategia de muestreo.
\end{knowledgebox}

\subsection{De la investigación a la producción: lista de verificación del lanzamiento gradual}
\begin{warningbox}{Las cuatro verificaciones que deben completarse antes de salir a producción}
\begin{itemize}
    \item \textbf{Lanzamiento gradual de corrección}: limitar primero el tráfico en tareas de alto valor, para verificar la tasa de aprobación real;
    \item \textbf{Lanzamiento gradual de costo}: comparar el beneficio/costo de la decodificación única frente al test-time scaling;
    \item \textbf{Lanzamiento gradual de estabilidad}: observar si las salidas de cadena larga presentan deriva de formato o razonamiento repetido;
    \item \textbf{Lanzamiento gradual de seguridad}: realizar pruebas de regresión de equipo rojo sobre prompts de alto riesgo y solicitudes que exceden los permisos autorizados.
\end{itemize}
\end{warningbox}

\subsection{Resumen de la sección}
Que la receta de entrenamiento genere valor a largo plazo depende de si se institucionaliza como un proceso del equipo. Un ritmo fijo, un registro completo y las verificaciones del lanzamiento gradual son lo que permite convertir los resultados de investigación en una capacidad de producto estable.

